%% paper/sections/conclusion.tex
%% Draft only. Brief, consistent with the all-null framing throughout.
\label{sec:conclusion}

This paper tested whether giving a video transformer tokens with a genuine
persistent referent changes what it learns, using a control designed to
isolate persistence from the entity-localized cropping and reduced token
budget that ordinarily travel with it. Under a decision rule and effect-size
floor frozen before any test-split evaluation, all three pre-registered
structural hypotheses came back null: MELD's persistence and entity-structure
margins are positive with confidence intervals excluding zero but short of
the pre-registered floor, and the entity-structure comparison's replication
on DAiSEE is indistinguishable from zero outright
(\S\ref{sec:results}--\ref{sec:discussion}). We report this as the finding,
not as a setback to be explained away: a pre-registered null, arrived at
through a control built specifically to make the positive result easy to
see if it were there, is evidence about the hypothesis, not an absence of
result.

Separately, and unconditionally on how the hypotheses resolved, entity-
localized cropping cuts per-crop encoder cost by $4.53\times$, the attention
mechanism this architecture is built around accounts for a fraction of a
percent of total compute, and the end-to-end cost of trading a spatial grid
for tracked entities depends on where detection is run --- a
$0.23\times$--$7.49\times$ range, not a single favorable number
(\S\ref{sec:efficiency}). Together with the dataset-admission procedure this
work produced as a reusable artifact (\S\ref{sec:dataset-audit}), the
contribution of this paper is less a positive architectural result than a
demonstration of what a fully pre-registered, matched-feature, cost-honest
evaluation of an intuitive idea actually looks like when the idea does not
clear its own bar.
